\section{Conclusion}
\label{sec:conclusion}

We presented a few-shot continual learning framework for 3D brain MRI using frozen foundation models and task-specific LoRA adapters. In the continual T2$\to$T3 setting (tumor segmentation then brain age regression, n\_shot=32), sequential full fine-tuning suffers severe catastrophic forgetting (T1 Dice 0.80$\to$0.16, BWT$\approx$-0.65), while sequential linear probing achieves strong T1 (0.79) but fails on T2 regression (MAE 1.45). EWC attains the lowest T2 MAE (0.001) among continual methods but exhibits severe forgetting (BWT$\approx$-0.65); LwF and Replay achieve strong T1 Dice ($\sim$0.79--0.80) yet suffer moderate to severe forgetting (BWT$\approx$-0.56 to $-0.78$). Our LoRA approach achieves the best balanced performance: competitive T2 MAE (0.012$\pm$0.003) among continual methods with \textbf{zero forgetting} (BWT=0) and $<$0.1\% trainable parameters per task. Encoder+decoder LoRA is necessary for segmentation; encoder-only yields poor Dice (0.19). LoRA is the only method that simultaneously maintains competitive performance on both tasks without catastrophic forgetting. Phase 3 (t32) and n\_shot=64 results are in Tables~\ref{tab:t32} and~\ref{tab:shot64}.
